%!TEX root = ../main.tex

\chapter{Evaluation Metrics\label{chap:evaluation}}

Evaluating the performance of a multi-class classification model for journalistic article quality requires careful selection of metrics, especially given the class imbalance in the dataset. As described in the Section~\ref{sec:dataset_statistics}, the four quality categories are not equally represented. This severe imbalance means that naive models could achieve seemingly high performance by favoring the majority class. Evaluation metrics must be chosen to reflect performance across all quality levels, rather than be dominated by the majority class.

\section{Limitations of Accuracy}
Accuracy measures the proportion of correct predictions out of all predictions. While it is a simple and intuitive metric, accuracy is inadequate for imbalanced datasets. In our context of quality classification, accuracy treats all prediction errors uniformly and thus can be misleading. A model labeling every article as \ac{Class 1} would achieve 54\% accuracy by virtue of Class 1’s prevalence, potentially outperforming more nuanced models that try to distinguish all four classes but incur some mistakes on the majority class. This illustrates a key limitation: \textbf{accuracy can be artificially high even if the model entirely ignores minority classes.} 

For a journalistic quality assessment system, such behavior is unacceptable – correctly identifying the minority classes (e.g. truly high-quality pieces) is crucial. Therefore, relying on accuracy alone would not highlight a model’s failure to detect rare yet important cases. In light of this, we turn to alternative metrics that better capture performance across all classes, giving due weight to both common and uncommon article categories. These metrics focus on the balance between types of errors, rather than just the overall fraction correct.

\section{Precision}

\textbf{Precision} is calculated as:
\begin{equation}
\text{Precision} = \frac{TP}{TP + FP},
\end{equation}
where \ac{TP} is the number of correctly predicted instances of a class, and \ac{FP} is the number of incorrectly predicted instances where the model predicted the positive class but the actual label was negative.

Precision answers the question: “\textit{When the model predicts a certain class (especially a positive or high-quality outcome), how often is it correct?}” In the context of our highest class (\ac{Class 3}), a high precision means that whenever the model flags an article as high-quality, it is usually correct. This is crucial for a platform like Seznam.cz because high precision in identifying top-quality content ensures that any article highlighted or promoted as “high-quality” truly meets the standards.

\section{Recall}

\textbf{Recall} measures the fraction of actual positive instances correctly identified by the model, calculated as:
\begin{equation}
\text{Recall} = \frac{TP}{TP + FN},
\end{equation}
where \ac{FN} is the number of actual positive instances that were incorrectly predicted as negative. It addresses “\textit{How many of the actual instances of a class did the model successfully find?}” For \ac{Class 3}, high recall means the model manages to identify most of the truly high-quality articles in the dataset.

\section{F1 Score}
The balance between precision and recall is a classic trade-off: improving one can often degrade the other. In our case, both are important. To illustrate:

\textbf{High Precision} – Ensures that articles labeled as high-quality truly meet quality standards, avoiding the risk of overestimating less suitable content.

\textbf{High Recall} – Ensures that most high-quality articles are identified correctly, maximizing the coverage of valuable content.

Depending on the application, one might prioritize precision over recall or vice versa. For instance, if the goal is to minimize the chance of misclassifying lower-quality content as top-tier, precision would take precedence. Alternatively, if the aim is to capture as many high-quality articles as possible, recall becomes more important—even at the cost of occasional misclassifications. In this thesis, we aim for a balance, since both types of errors carry costs.

To evaluate model performance holistically, we employ the \textbf{F1 score}, the harmonic mean of precision and recall. The F1 score yields a high value only when both precision and recall are high, making it well-suited for our imbalanced multi-class problem. It is defined as:

\begin{equation}
\text{F1} = 2 \times \frac{\text{Precision} \times \text{Recall}}{\text{Precision} + \text{Recall}}
\end{equation}

This formulation captures the trade-off between \ac{FP} and \ac{FN}, ensuring that the model cannot achieve a high score by excelling in only one aspect.

Unlike accuracy—which treats all errors equally and can be skewed by majority classes—the F1 score ensures that minority classes are not overshadowed. A high F1 means the model is correctly identifying most high-quality articles (high recall) while avoiding mislabeling low-quality ones (high precision). F1 was used as the primary metric during model development.

\section{Weighted F1 Score}
In multi-class classification, extending precision, recall, and F1 requires aggregation across classes. For this project, the \textbf{weighted F1 score} is used as the main evaluation metric, accounting for varying class frequencies. It is calculated by computing the F1 score for each class and taking their weighted average using class support (i.e., number of instances) as weights:

\begin{equation}
  \text{F1W} = \sum_{i=1}^{n} (w_i \times \text{F1}_i)
  \end{equation}
where $i \in \{0,1,2,3\}$, \(w_i = \frac{N_i}{N}\) is the fraction of samples belonging to class \(i\) in the dataset (i.e. its empirical probability) and \(\text{F1}_i\) is the F1 score for class \(i\).
This gives more influence to more common classes while still including all categories.

\subsection{Why use weighted F1?}
In imbalanced multi-class tasks like this one, the weighted F1 offers key benefits over macro or micro averaging:

\begin{itemize}
\item \textbf{Balances Class Influence:} Macro-averaged F1 weighs all classes equally, possibly overstating rare class performance. Micro-averaging, on the other hand, heavily reflects majority class behavior. Weighted F1 uses class prevalence to moderate both extremes.
\item \textbf{Comprehensive Single Metric:} Weighted F1 produces a single summary that reflects performance across all classes, aiding model selection and comparison. During development, this metric was used on the validation set to guide hyperparameter tuning and model selection (\ref{sec:training_methodology}).
\item \textbf{Reflects Real Data:} It mirrors the actual class distribution in the data, ensuring that both common and rare classes are appropriately considered without letting any one class disproportionately affect evaluation.
\end{itemize}

\section{Confusion Matrix}
The confusion matrix is a valuable visualization for understanding classification performance. It shows how instances of each true class are predicted across all classes, helping reveal systematic misclassifications. See \reffig{fig:confusion_example} for an example.

\begin{figure}
\centering
\includegraphics[width=0.8\textwidth]{./fig/conf_matrix/confusion_matrix_Deberta_large.png}
\caption{Example of Confusion Matrix. DebertaV3 Large}
\label{fig:confusion_example}
\end{figure}

In this study, the weighted F1 score was selected as the most appropriate evaluation metric and serves as the primary basis for reporting results. It was also integral to training: the final model was chosen based on its highest validation weighted F1, and early stopping was applied once this score plateaued. Using this metric consistently ensures alignment between training and evaluation phases, even under class imbalance.

\section{Business Relevance}

The choice of evaluation metrics extends beyond technical considerations; it reflects practical goals for platforms that aim to assess article quality automatically. The model is designed to approximate human editors' evaluations across a quality spectrum from \ac{Class 0} to \ac{Class 3}. As discussed in the Chapter \ref{chap:methodology}, human-provided quality labels set a performance ceiling, and evaluation metrics should reflect how closely the model approaches that benchmark.

The weighted F1 score supports this goal by measuring model effectiveness across all quality levels, proportionally to their occurrence in the dataset. A high weighted F1 indicates that the model correctly distinguishes between different article quality levels in a way that mirrors editorial standards. This helps ensure consistency in automated quality assessment—recognizing stronger pieces and identifying weaker ones accurately, without favoring one class excessively. Ultimately, this facilitates editorial scalability while maintaining reliable quality judgments.

\subsection{Cost of Classification Errors}
Different misclassification types affect the system differently:

\begin{itemize}
\item \textbf{False Positives (FPs):} When lower-quality articles are mistakenly marked as high-quality, the perceived reliability of recommendations may decline. Users could become skeptical of the platform's ability to surface genuinely strong content.
\item \textbf{False Negatives (FNs):} When high-quality articles are not correctly identified, potential value is lost, as outstanding content may fail to reach its intended audience, lowering engagement and satisfaction.
\end{itemize}

The F1 score addresses both risks by balancing precision and recall. Its weighted variant further ensures that infrequent but important categories are not neglected. By incorporating these considerations into the evaluation process, the model better supports content selection goals, ensuring that all quality levels are appropriately represented in automatic assessments.
